\documentclass[10pt,a4paper]{amsart}
\usepackage[margin=19mm]{geometry}
\usepackage{amsmath,amssymb,mathtools}
\usepackage[scr=rsfs,cal=euler]{mathalfa}
\usepackage{xfrac}
\usepackage[unicode,hidelinks,bookmarks=false]{hyperref}
\newcommand{\ie}{i.e.\ }
\newcommand{\cf}{cf.\ }
\newcommand{\resp}{respectively\ }
\newcommand{\Subsection}{\subsection}
\providecommand{\nobreakdash}{\nobreak\hbox{-}}
\setlength{\emergencystretch}{2em}
\multlinegap0pt
\usepackage{amssymb}
\usepackage{mathtools}
\usepackage[shortlabels]{enumitem}
\usepackage{xcolor}
\usepackage{bm}
\usepackage[normalem]{ulem}
\usepackage{tikz}
\newtheorem{proposition}{Proposition}
\newtheorem{lemma}{Lemma}
\def\cleq{\preccurlyeq}
\title{Smaller universal posets}
\author{Paul Bastide, Carla Groenland, Rajko Nenadov}
\date{Source: arXiv:2509.17820v3 (CC BY 4.0)}
\begin{document}
\maketitle

\noindent\emph{Source and licence.} Smaller universal posets. Version arXiv:2509.17820v3 (CC BY 4.0), distributed by arXiv under the Creative Commons Attribution 4.0 International licence, http://creativecommons.org/licenses/by/4.0/, as recorded in the arXiv licence metadata for this exact version and shown on the abstract page https://arxiv.org/abs/2509.17820v3. Full licence notice: \texttt{licenses/arxiv-2509.17820-CC-BY-4.0.txt}.\par\smallskip

\noindent\textit{Editorial scope.} Counted unit: the article's lemma lem:hypercube\_antichain (source0.tex lines 370--374). The article's complete original proof of it is delivered with it (source0.tex lines 376--432, one \texttt{proof} environment of 57 lines, which the article places away from the statement). No mathematics is written, repaired or re-derived: the statement and the proof are the article's own lines, in the article's own notation, and no retained line is substituted or re-flowed.  Retained uncounted as the source-local context the proof needs: lines 301--304.  Nothing was omitted from any retained span, so the package carries no omission record.  No retained line contains a citation, so the wrapper carries no bibliography and no bibliography entry.  No retained line contains a figure environment, an image-inclusion command or drawing code, so no image is delivered and the package declares no asset dependency.  Editorial formatting only, in addition: an AMS \texttt{amsart} preamble that restates the article's own declarations and macros used by the retained lines, namely the environments \texttt{proposition}, \texttt{lemma} on separate counters, together with the standard packages those lines need.  The retained environments print in this document as Proposition 1, Lemma 1 in order of appearance, whereas the article prints the same items with the numbering of its own full document.  The complete native source of the article as distributed in the arXiv e-print for this exact version is bundled unchanged as \texttt{sources/185481/source0.tex}.
\begin{proposition}
    \label{prop:hypercube_univ}
    The poset $(2^{[n]},\subseteq)$ is universal for the class of all $n$-element posets.
\end{proposition}

\begin{lemma}
\label{lem:hypercube_antichain}
    Let $n\geq a\geq 2$ be integers. Let $\mathcal{A}_n(a)$ denote the family of all $n$-element posets $P$ which contain an antichain of size $a$. Let $\ell$ be the smallest integer for which $\binom{\ell}{\lfloor\ell/2\rfloor}\geq a$. 
    Then the Boolean lattice $(2^{[n-a+\ell]},\subseteq)$ is a universal poset for $\mathcal{A}_n(a)$. 
\end{lemma}

\begin{proof}[Proof of Lemma~\ref{lem:hypercube_antichain}]
In this proof, we use the notation $[a,b]$ for the set of integers $x$ with $a\leq x\leq b$.

    Let $(P,\cleq_P) \in \mathcal{A}_n(a)$ and let $A \subseteq P$ be an antichain of size $a$. Let $b$ denote the number of elements $x\in P$ with $x\cleq_P y$ for some $y\in A$. We number all such elements of $P$ ``below'' $A$ by $x_1,\dots,x_b$. We assign each element in the antichain $A$ a unique (and arbitrary) subset $S\subseteq [b+1,b+\ell]$ of size $\lfloor \ell/2 \rfloor$ and name the element itself $y_S$. (Note that this is possible by our choice of $\ell$.) We number the remaining $n-a-b$ elements by $z_{b+\ell+1},\dots,z_{n-a+\ell}$.  

    The elements $x_i$ can be seen as ``below'' the antichain and the elements $z_j$ as ``above'' the antichain in the following sense. Firstly, $y_S\not\cleq_P x_i$ for each $i$ and $S$ (using that $A$ is an antichain, we cannot have $y_{S}\cleq_P x_i\cleq_P y_{S'}$). Secondly, $z_j\not\cleq_P y_S$ and $z_j\not\cleq_P x_i$ as $z_j$ is not ``below'' an element in the antichain by definition and $x_i$ is. 
    Given $v \in P$, we define $f(v)\in 2^{[n-a+\ell]}$ depending on whether $v$ is ``below'' $A$, in $A$ or ``above'' $A$.

    For $i\in [1,b]$, 
    \[
    f(x_i)=\{j\in [b]:x_j\preceq x_i\}.
    \]
For $S\subseteq [b+1,b+\ell]$, if $y_S$ is defined, \[
f(y_S)=\{i\in [b]:x_i\cleq_P y_S\}\cup S\cup  \{j\in [b+\ell+1,n-a+\ell]:y_S\not \cleq_P z_j\}.
\]
For $j\in [b+\ell+1,n-a+\ell]$,  \[
f(z_j)=\{i\in [b]:x_i\cleq_P z_j\} \cup [b+1,b+\ell]
\cup 
\{i\in [b+\ell+1,n-a+\ell]:z_j \not \cleq_P z_i\}. 
\]
What remains to show is that for all elements $u,v\in P$, we have
\begin{equation}
    \label{eq:poset_contains}
u\cleq_P v \iff f(u)\subseteq f(v).
\end{equation}

\paragraph{Case 1: $u = x_i$.} Note that, similar to the proof of  Proposition~\ref{prop:hypercube_univ}, for all $i\in [b]$ and $v\in P$ we have $x_i\cleq_P v$ if and only if $i\in f(v)$. This implies that
\[
x_i\cleq_P v \iff f(x_i)\subseteq f(v)
\]
for the same reason as in that proof.
Indeed, if  $x_i\not\cleq_Pv$, then $i\in f(x_i)\setminus f(v)$. On the other hand, if $x_i\cleq_P v$, then for all $j\in f(x_i)$, $x_j\cleq_P x_i$ implies that $x_j\cleq_Pv$  and so $j\in f(v)$. 

\paragraph{Case 2: $u = y_S$.} 
We consider the case $u=y_S\in A$ by case analyis on $v$.
\begin{itemize}
    \item If $v=x_i$ for some $i\in [1,b]$, then $y_S\not\cleq_P x_i$ by the definition of ``below''. We also find that $S\subseteq f(y_S)\setminus f(x_i)$ and so $f(y_S)\not\subseteq f(x_i)$, as desired.
    \item If $v=y_{S'}\in A$ with $u\neq v$, then $u,v$ are incomparable in $P$. As $S$ and $S'$ have the same size and $S \neq S'$, we conclude $S \not \subseteq S'$. Therefore, $f(u) \not \subseteq f(v)$.
    
    \item If $v=z_j$ for some $j\in [b+\ell+1, n - a + \ell]$ and $y_S\not\cleq z_j$, we find $j\in f(y_S)\setminus f(z_j)$, where $j \not \in f(z_j)$ since $z_j \cleq_P z_j$. So $f(y_s)\not \subseteq f(z_j)$. 
    
    On the other hand, if $y_S\cleq_P z_j$, then for all $i\in [b]$ by definition of $f$ and transitivity \[
    i\in f(y_S)\implies 
    x_i\cleq_P y_S\implies x_i\cleq_P z_j \implies i\in f(z_j).
    \]
    For the same reasons for all $i\in [b+\ell+1,n-a+\ell]$,
    \[
    i\in f(y_S)\implies y_S\not\cleq_P z_i\implies z_j \not\cleq_P z_i\implies i\in f(z_j).
    \]
    By definition, $[b+1,b+\ell]\subseteq f(z_j)$. This shows that $f(y_S)\subseteq f(z_j)$ if $y_S\cleq_P z_j$.
\end{itemize}

\paragraph{Case 3: $u = z_j$.} Finally, we consider the case where $u=z_j$ for $j\in [b+\ell+1,n-a+\ell]$. Now $u\not\cleq v$ for all $v$ of the form $x_i$ and $y_S$ (by the definition of ``below'') and for those $v$ we also find that $f(u)\not\subseteq f(v)$ since $[b+1,b+\ell]\not\subseteq f(v)$ whereas $[b+1,b+\ell]\subseteq f(u)$. 
When $v=z_{j'}$ for $j'\in [b+\ell+1,n-a+\ell]$, the proof is analogous to before: if $z_j\not\cleq_P z_{j'}$, then $j'\in f(z_{j})\setminus f(z_{j'})$, whereas transitivity of $\cleq_P$ ensures that $f(z_j)\subseteq f(z_{j'})$ when $z_j\cleq_P z_{j'}$.

Since we proved (\ref{eq:poset_contains}) in all cases for $u$ and $v$, this shows that $(2^{[n-a+\ell]},\subseteq )$ contains the poset $P$ via the embedding given by $f$.
\end{proof}

\end{document}
